\documentclass[11pt,a4paper]{article}

\usepackage[utf8]{inputenc}
\usepackage[T1]{fontenc}
\usepackage{amsmath,amssymb}
\usepackage{enumitem}
\usepackage{booktabs}
\usepackage{geometry}
\usepackage{hyperref}
\usepackage{microtype}
\usepackage{parskip}

\geometry{margin=2.5cm}

\title{\textbf{Peer Review Report}\\[0.4em]
\large Frequency maximization of planar structures using quasi-static\\
approximation in topology optimization}
\author{Multi-model consensus review (GPT + Claude Sonnet)\\
\small Generated: 2026-05-22}
\date{}

\begin{document}

\maketitle
\thispagestyle{empty}

\hrule
\vspace{0.5em}
\noindent\textbf{Manuscript:} Frequency maximization of planar structures using
quasi-static approximation in topology optimization\\
\noindent\textbf{Authors:} P.\ Tauzowski, A.\ Szczepańczyk, B.\ Błachowski\\
\noindent\textbf{Review type:} Multi-model consensus (GPT + Claude Sonnet, independent pipelines)\\
\noindent\textbf{Review confidence:} HIGH (both models agree; no conflicting signals)
\vspace{0.5em}
\hrule

%-----------------------------------------------------------------------
\section*{Recommendation}
%-----------------------------------------------------------------------

\textbf{MAJOR REVISION}

The manuscript presents a genuine and useful computational contribution ---
a quasi-static approximation that reduces frequency-maximization topology
optimization to a standard compliance code after a single initial
eigensolve. The core algorithmic idea is coherent and the reported
computational savings at large mesh sizes are real. However, several
issues must be resolved before the paper is suitable for publication:
a data inconsistency in Table~3 directly contradicts a stated principal
contribution; the central approximation is not analytically or
numerically validated; and the benchmark comparison involves a
non-canonical comparator implementation that is not consistently
disclosed.

%-----------------------------------------------------------------------
\section{Conceptual and Theoretical Issues}
%-----------------------------------------------------------------------

\begin{enumerate}[leftmargin=*, label=\textbf{C\arabic*.}]

\item \textbf{Omitted load sensitivity --- central approximation unvalidated.}
Equation~(5) defines the inertial load $f(\mathbf{x}) =
\omega_1^{(0)2}\mathbf{M}(\mathbf{x})\boldsymbol{\Phi}_1$, which depends on
element densities through the evolving mass matrix $\mathbf{M}(\mathbf{x})$.
Equation~(6) drops the corresponding load-sensitivity term
$\partial f / \partial x_e = \omega_1^{(0)2}(\partial \mathbf{M}/\partial x_e)
\boldsymbol{\Phi}_1$ from the compliance gradient.
This omission is the defining simplification of the method, but it is
not quantitatively justified. The citation of Munro~(2024) in
Section~3.2 concerns the self-weight OC update, not the frozen-mode
sensitivity approximation, and does not establish the validity of this
specific approximation for the problem class studied.
A numerical ablation comparing Eq.~(6) against the full
design-dependent-load derivative, or against a periodically
updated-mode variant, is required on at least one benchmark.

\item \textbf{Multi-mode sensitivity derivation has no theoretical support.}
Section~3.3 combines two inertial loads via a scalar weight
$\alpha \in [0,1]$ (Eq.~7) and uses the same stiffness-only
sensitivity (Eq.~6) for the combined load, omitting load-sensitivity
contributions from both eigenvectors simultaneously.
No argument is provided for why omitting both terms jointly preserves a
useful descent direction for either targeted eigenfrequency, or for what
range of $\alpha$ the joint omission remains acceptable.

\item \textbf{Monotone gain scaling --- stated as principal contribution,
contradicted by Table~3.}
The abstract, Section~4.2, and the conclusions state that
``load-weighted modal gains scale monotonically with the assigned load
weight'' (bidirectional). For the $\widetilde{\omega}_2/\omega_2^{(0)}$
tracked-gain sequence in Table~3 (clamped beam), the values as
$\alpha$ decreases from 1.0 to 0 are: $1.99\times$, $1.73\times$,
$2.05\times$, $2.25\times$, $2.45\times$. The gain at $\alpha=0.75$
$(1.73\times)$ is lower than at $\alpha=1.0$ $(1.99\times)$, which
contradicts strict monotone decrease with decreasing $\alpha$.
Note: the building-example $\Phi_2$ gain sequence in Table~5 is
monotone. The failure is specific to the clamped-beam case at
$\alpha=0.75$. The claim must be verified and corrected, or qualified.

\item \textbf{Frozen-mode reliability not demonstrated.}
The conclusions acknowledge an ``accuracy ceiling'' from the frozen
initial mode shape, but the paper does not demonstrate conditions under
which the approximation fails. A reliability diagnostic --- for example,
a benchmark where the initial and optimized mode shapes differ
substantially, a MAC-based threshold analysis, or a comparison against
a periodically updated-mode variant --- is needed to bound the practical
scope of the method.

\end{enumerate}

%-----------------------------------------------------------------------
\section{Algorithmic and Numerical Methodology Concerns}
%-----------------------------------------------------------------------

\begin{enumerate}[leftmargin=*, label=\textbf{A\arabic*.}]

\item \textbf{The Olhoff comparator is a non-canonical in-house implementation.}
Section~5.1 identifies three modifications to the original Du--Olhoff
(2007)/Olhoff--Du (2014) procedure: (a)~a Heaviside projection with
seven-level continuation and a grayness penalty; (b)~the inner MMA
convergence loop collapsed to a single call per outer iteration; and
(c)~a trial eigensolve rejection step that ``effectively doubles the
number of eigensolves per outer iteration.'' These deviations
collectively increase the comparator's runtime relative to the
canonical formulation. The 7.1$\times$ speedup and 8.6\% frequency gap
must be presented as comparisons against this specific in-house variant.
A compact table listing what each implementation (original Olhoff--Du,
in-house comparator, Yuksel--Yilmaz, proposed method) includes or
excludes would clarify the comparison.

\item \textbf{Inconsistent optimizer within the proposed method.}
Algorithm~3.5 and Section~3.4 describe an OC bisection update.
Sections~4.2 and~4.3 state that the MMA optimizer with move limit~0.2
is used. The proposed method therefore uses different optimizers across
its own benchmark cases, without discussion of whether this affects
iteration counts, convergence speed, or frequency results.

\item \textbf{Number of initialization modes unspecified.}
Step~2 of Algorithm~3.5 reads ``Solve eigenvalue problem
$\mathbf{K}_0\boldsymbol{\Phi}_j = \lambda_j\mathbf{M}_0\boldsymbol{\Phi}_j$,
$j=1,\ldots,n_\mathrm{modes}$'' but $n_\mathrm{modes}$ is never stated
as a parameter. Its value must be specified for each benchmark.

\item \textbf{Inconsistent void pseudo-material parameters.}
Section~4.1 (simply supported beam) uses $E_\mathrm{min}=10^{-9}E_0$;
Section~4.2 (clamped beam) uses $E_\mathrm{min}=10^{-6}E_0$. The
basis for this difference is not explained.

\end{enumerate}

%-----------------------------------------------------------------------
\section{Validation and Benchmarking Weaknesses}
%-----------------------------------------------------------------------

\begin{enumerate}[leftmargin=*, label=\textbf{V\arabic*.}]

\item \textbf{Table~1 lacks statistical uncertainty.}
Runtime and memory values are labeled ``averaged over 10 independent
runs'' without standard deviations. At the $160\times20$ mesh, the
proposed method (5.1~s) is marginally slower than the Yuksel code
(4.7~s). Without variance information, this difference and others in
Table~1 cannot be interpreted statistically. Standard deviations (or
at minimum, ranges) must be reported; the source of run-to-run
variability in a deterministic MATLAB code should be explained.

\item \textbf{No eigenfrequency values across mesh resolutions.}
Table~1 reports runtime and iterations at four mesh sizes but does not
include the achieved $\omega_1$ for any method at any resolution.
The 8.6\% frequency gap is stated only for the $400\times50$ case.
A column for $\omega_1$ at each mesh size should be added so that the
mesh-dependence of the frequency gap can be assessed.

\item \textbf{No mesh independence for multi-mode benchmarks.}
The clamped-beam ($400\times50$) and building ($80\times240$) examples
are each tested at a single mesh. MAC-tracked gains, mode-ordering
shifts, and the $\alpha$-behavior conclusions should be verified at
at least one additional discretization.

\item \textbf{Low MAC$_2$ values in Table~5 not acknowledged.}
Table~5 reports $\mathrm{MAC}_2=0.520$ at $\alpha=0.25$ and
$\mathrm{MAC}_2=0.635$ at $\alpha=0.5$ for the building example.
These values indicate that the mode identified as best matching
$\boldsymbol{\Phi}_2$ shares only about 27\% and 40\% of its modal
energy with the reference, respectively. The frequency gains reported
for these configurations (3.32$\times$ and 3.09$\times$) are based on
poorly correlated modes. The authors should define a MAC validity
threshold and qualify or exclude gain figures below it.

\item \textbf{Initial clamped-beam eigenfrequencies not stated.}
Section~4.3 explicitly states $\omega_1^{(0)}=19.84$~rad/s and
$\omega_2^{(0)}=90.93$~rad/s for the building domain. The analogous
values for the clamped-beam domain are absent from Section~4.2.
Table~3 gain factors cannot be independently verified without them.

\item \textbf{Headline claims must be scoped.}
\begin{itemize}
  \item ``within 8.6\% of the bound-formulation optimum'' should be
    qualified to the simply supported beam benchmark at $400\times50$.
  \item ``completing the optimization 7.1$\times$ faster'' is specific to
    $400\times50$; at $160\times20$ the proposed method is marginally
    slower than the Yuksel code.
  \item ``reaching up to 4.61$\times$ over the initial solid-domain
    frequency'' refers to the MAC-tracked $\Phi_1$-type gain, not the
    conventional sorted first eigenfrequency. This distinction should
    appear in the abstract.
\end{itemize}

\end{enumerate}

%-----------------------------------------------------------------------
\section{Mathematical Notation and Rigor}
%-----------------------------------------------------------------------

\begin{enumerate}[leftmargin=*, label=\textbf{M\arabic*.}]

\item \textbf{Eigenvector notation inconsistency.}
The manuscript uses uppercase $\boldsymbol{\Phi}$ (Eqs.~5, 7,
Section~3.6) and lowercase $\boldsymbol{\phi}$ (Eqs.~2, 3, MAC
definition Eq.~9) for eigenvectors without a stated convention.
Equation~(9) uses both $\boldsymbol{\phi}_i$ and $\boldsymbol{\Phi}_j$
in the same expression, suggesting an intentional distinction (current
vs.\ reference eigenvector), but this is never defined. A convention
should be stated at first use and applied consistently.

\item \textbf{Sensitivity filter notation (Eq.~8).}
The filtered sensitivity symbol is introduced without a formal
definition. The denominator summation domain ($\sum H_{ei}$) should
state explicitly that the sum runs over the neighborhood $N_e$ of
element~$e$.

\item \textbf{Cross-table discrepancy at $\alpha=1.0$ (clamped beam).}
Table~2 gives $\omega_2=365.78$~rad/s; Table~3 gives the
$\Phi_1$-tracked mode (MAC$=0.995$) at topo-mode~3 with
$\widetilde{\omega}_1=368.2$~rad/s. The 2.4~rad/s difference
($\approx0.7\%$) between these quantities should be explained
(rounding, independent runs, or precision issue).

\item \textbf{Scaling claim not supported by reported data.}
Section~5.2 states ``both methods exhibit an empirical per-iteration
scaling of roughly $\mathcal{O}(n_e^{1.3})$.'' The per-iteration times
in Table~1 are consistent with exponents of approximately 1.06 (Olhoff)
and 0.8 (Yuksel) over the four reported mesh sizes, substantially
below~1.3. The claim should be supported with a log-log scaling figure
or the specific exponent values should be corrected.

\end{enumerate}

%-----------------------------------------------------------------------
\section{Reproducibility and Transparency}
%-----------------------------------------------------------------------

\begin{enumerate}[leftmargin=*, label=\textbf{R\arabic*.}]

\item \textbf{No versioned code archive.}
The GitHub URL is provided without a commit hash or archived release
DOI. Repository contents may change after publication. An immutable
archive (e.g., Zenodo DOI or GitHub release tag) should be provided.

\item \textbf{Hardware and environment not specified.}
CPU model, clock speed, RAM, operating system, MATLAB version, and
thread count are absent. These are necessary to assess the reported
timing figures.

\item \textbf{No script-to-figure mapping.}
The manuscript does not specify which script or configuration file
reproduces each figure or table entry.

\item \textbf{``No eigensolve machinery'' is inaccurate.}
The abstract states the method requires ``only a single eigensolve at
initialization.'' Algorithm~3.5 Step~5 performs an additional eigensolve
at convergence. The correct description is ``no eigensolve \emph{inside}
the optimization loop.''

\end{enumerate}

%-----------------------------------------------------------------------
\section{Required Revisions}
%-----------------------------------------------------------------------

\subsection*{Critical (must be resolved before reconsideration)}

\begin{enumerate}[leftmargin=*, label=\textbf{CR\arabic*.}]

\item Verify and correct Table~3 $\Phi_2$-tracked gain data ($\alpha=0.75$
  entry shows $1.73\times$, which is below the $\alpha=1.0$ entry of
  $1.99\times$). Provide diagnostic output (iteration history, MAC
  evolution) for the $\alpha=0.75$ clamped-beam case. Revise or qualify
  all monotonicity claims in the abstract, contribution list,
  Section~4.2, and conclusions to match the corrected or explained data.

\item Add a quantitative validation of the omitted load sensitivity
  (Eq.~6): compare the approximate gradient against the full
  design-dependent-load derivative, or against a periodically
  updated-mode variant, on at least one benchmark.

\item Relabel the Olhoff comparator as an in-house Olhoff-inspired
  implementation throughout the manuscript; provide a comparison table
  showing what each method includes or excludes; state quantitatively
  how the three deviations (Section~5.1) contribute to the reported
  speedup and frequency gap.

\end{enumerate}

\subsection*{Major (must be addressed)}

\begin{enumerate}[leftmargin=*, label=\textbf{MR\arabic*.}]

\item Add standard deviations to all runtime and memory entries in
  Table~1; explain the source of variability across the 10~runs; specify
  CPU, RAM, OS, MATLAB version, and thread count.

\item Add eigenfrequency ($\omega_1$) values for all three methods at all
  four mesh resolutions in Table~1 or a supplementary table.

\item State initial clamped-beam eigenfrequencies $\omega_1^{(0)}$ and
  $\omega_2^{(0)}$ in Section~4.2.

\item Define and apply a MAC validity threshold; qualify or exclude frequency
  gains for configurations below the threshold (particularly
  $\alpha=0.25$ and $\alpha=0.5$ in Table~5).

\item Provide a mesh convergence study (minimum two mesh sizes) for the
  clamped-beam or building multi-mode benchmarks.

\item Provide a frozen-mode reliability diagnostic: at minimum, one
  benchmark or analysis characterizing when the initial mode shape is
  a poor proxy for the optimal mode shape, or demonstrating that the
  three tested problem classes lie within the reliable regime.

\item Scope all headline quantitative claims to the specific benchmark,
  mesh, and metric to which they apply.

\end{enumerate}

\subsection*{Minor (should be corrected)}

\begin{enumerate}[leftmargin=*, label=\textbf{mn\arabic*.}]

\item Correct figure references in Section~4.1 (``Figs.~5--7'' should
  be ``Figs.~1--3'').
\item Define the eigenvector notation convention ($\boldsymbol{\Phi}$
  vs.\ $\boldsymbol{\phi}$) at first use and apply consistently.
\item Complete the Olhoff--Du (2014) reference (publisher location listed
  as ``???'').
\item Archive the code at a permanent DOI and report the commit hash.
\item Specify hardware and MATLAB configuration in Section~5 or a
  reproducibility statement.
\item Provide a script-to-figure mapping in the repository README.
\item Address the OC-vs-MMA optimizer switch: state which optimizer is
  used for the proposed method in each benchmark and justify the choice.
\item Correct or support the $\mathcal{O}(n_e^{1.3})$ scaling claim
  (Section~5.2) with a log-log plot or revised exponent values.
\item Justify or unify the void pseudo-material parameter differences
  ($E_\mathrm{min}=10^{-9}E_0$ vs.\ $10^{-6}E_0$) across benchmarks.
\item Correct the inaccurate ``no eigensolve machinery'' claim to
  ``no eigensolve inside the optimization loop.''
\item Clarify the 2.4~rad/s discrepancy between Tables~2 and~3 at
  $\alpha=1.0$ (clamped beam, $\omega_2=365.78$ vs.\
  $\widetilde{\omega}_1=368.2$~rad/s).

\end{enumerate}

%-----------------------------------------------------------------------
\section*{Cross-Review Note}
%-----------------------------------------------------------------------

This report is based on an independent dual-model review (GPT and Claude
Sonnet), with both models reaching MAJOR REVISION with HIGH confidence.
The decision and trusted issue list are derived from the consensus of
both models; issues raised by only one model are retained where they
are supported by manuscript-visible evidence and are marked accordingly
in the editorial record. No hallucinated or invalid reviewer comments
were identified in either model's pipeline. The consensus weighted score
is approximately 2.80 on a 1--5 scale; the MAJOR REVISION recommendation
reflects an editorial override based on the correctable nature of the
critical issues.

\end{document}
\documentclass[11pt,a4paper]{article}

\usepackage[utf8]{inputenc}
\usepackage[T1]{fontenc}
\usepackage{amsmath,amssymb}
\usepackage{enumitem}
\usepackage{booktabs}
\usepackage{geometry}
\usepackage{hyperref}
\usepackage{microtype}
\usepackage{parskip}

\geometry{margin=2.5cm}

\title{\textbf{Peer Review Report}\\[0.4em]
\large Frequency maximization of planar structures using quasi-static\\
approximation in topology optimization}
\author{Multi-model consensus review (GPT + Claude Sonnet)\\
\small Generated: 2026-05-22}
\date{}

\begin{document}

\maketitle
\thispagestyle{empty}

\hrule
\vspace{0.5em}
\noindent\textbf{Manuscript:} Frequency maximization of planar structures using
quasi-static approximation in topology optimization\\
\noindent\textbf{Authors:} P.\ Tauzowski, A.\ Szczepańczyk, B.\ Błachowski\\
\noindent\textbf{Review type:} Multi-model consensus (GPT + Claude Sonnet, independent pipelines)\\
\noindent\textbf{Review confidence:} HIGH (both models agree; no conflicting signals)
\vspace{0.5em}
\hrule

%-----------------------------------------------------------------------
\section*{Recommendation}
%-----------------------------------------------------------------------

\textbf{MAJOR REVISION}

The manuscript presents a genuine and useful computational contribution ---
a quasi-static approximation that reduces frequency-maximization topology
optimization to a standard compliance code after a single initial
eigensolve. The core algorithmic idea is coherent and the reported
computational savings at large mesh sizes are real. However, several
issues must be resolved before the paper is suitable for publication:
a data inconsistency in Table~3 directly contradicts a stated principal
contribution; the central approximation is not analytically or
numerically validated; and the benchmark comparison involves a
non-canonical comparator implementation that is not consistently
disclosed.

%-----------------------------------------------------------------------
\section{Conceptual and Theoretical Issues}
%-----------------------------------------------------------------------

\begin{enumerate}[leftmargin=*, label=\textbf{C\arabic*.}]

\item \textbf{Omitted load sensitivity --- central approximation unvalidated.}
Equation~(5) defines the inertial load $f(\mathbf{x}) =
\omega_1^{(0)2}\mathbf{M}(\mathbf{x})\boldsymbol{\Phi}_1$, which depends on
element densities through the evolving mass matrix $\mathbf{M}(\mathbf{x})$.
Equation~(6) drops the corresponding load-sensitivity term
$\partial f / \partial x_e = \omega_1^{(0)2}(\partial \mathbf{M}/\partial x_e)
\boldsymbol{\Phi}_1$ from the compliance gradient.
This omission is the defining simplification of the method, but it is
not quantitatively justified. The citation of Munro~(2024) in
Section~3.2 concerns the self-weight OC update, not the frozen-mode
sensitivity approximation, and does not establish the validity of this
specific approximation for the problem class studied.
A numerical ablation comparing Eq.~(6) against the full
design-dependent-load derivative, or against a periodically
updated-mode variant, is required on at least one benchmark.

\item \textbf{Multi-mode sensitivity derivation has no theoretical support.}
Section~3.3 combines two inertial loads via a scalar weight
$\alpha \in [0,1]$ (Eq.~7) and uses the same stiffness-only
sensitivity (Eq.~6) for the combined load, omitting load-sensitivity
contributions from both eigenvectors simultaneously.
No argument is provided for why omitting both terms jointly preserves a
useful descent direction for either targeted eigenfrequency, or for what
range of $\alpha$ the joint omission remains acceptable.

\item \textbf{Monotone gain scaling --- stated as principal contribution,
contradicted by Table~3.}
The abstract, Section~4.2, and the conclusions state that
``load-weighted modal gains scale monotonically with the assigned load
weight'' (bidirectional). For the $\widetilde{\omega}_2/\omega_2^{(0)}$
tracked-gain sequence in Table~3 (clamped beam), the values as
$\alpha$ decreases from 1.0 to 0 are: $1.99\times$, $1.73\times$,
$2.05\times$, $2.25\times$, $2.45\times$. The gain at $\alpha=0.75$
$(1.73\times)$ is lower than at $\alpha=1.0$ $(1.99\times)$, which
contradicts strict monotone decrease with decreasing $\alpha$.
Note: the building-example $\Phi_2$ gain sequence in Table~5 is
monotone. The failure is specific to the clamped-beam case at
$\alpha=0.75$. The claim must be verified and corrected, or qualified.

\item \textbf{Frozen-mode reliability not demonstrated.}
The conclusions acknowledge an ``accuracy ceiling'' from the frozen
initial mode shape, but the paper does not demonstrate conditions under
which the approximation fails. A reliability diagnostic --- for example,
a benchmark where the initial and optimized mode shapes differ
substantially, a MAC-based threshold analysis, or a comparison against
a periodically updated-mode variant --- is needed to bound the practical
scope of the method.

\end{enumerate}

%-----------------------------------------------------------------------
\section{Algorithmic and Numerical Methodology Concerns}
%-----------------------------------------------------------------------

\begin{enumerate}[leftmargin=*, label=\textbf{A\arabic*.}]

\item \textbf{The Olhoff comparator is a non-canonical in-house implementation.}
Section~5.1 identifies three modifications to the original Du--Olhoff
(2007)/Olhoff--Du (2014) procedure: (a)~a Heaviside projection with
seven-level continuation and a grayness penalty; (b)~the inner MMA
convergence loop collapsed to a single call per outer iteration; and
(c)~a trial eigensolve rejection step that ``effectively doubles the
number of eigensolves per outer iteration.'' These deviations
collectively increase the comparator's runtime relative to the
canonical formulation. The 7.1$\times$ speedup and 8.6\% frequency gap
must be presented as comparisons against this specific in-house variant.
A compact table listing what each implementation (original Olhoff--Du,
in-house comparator, Yuksel--Yilmaz, proposed method) includes or
excludes would clarify the comparison.

\item \textbf{Inconsistent optimizer within the proposed method.}
Algorithm~3.5 and Section~3.4 describe an OC bisection update.
Sections~4.2 and~4.3 state that the MMA optimizer with move limit~0.2
is used. The proposed method therefore uses different optimizers across
its own benchmark cases, without discussion of whether this affects
iteration counts, convergence speed, or frequency results.

\item \textbf{Number of initialization modes unspecified.}
Step~2 of Algorithm~3.5 reads ``Solve eigenvalue problem
$\mathbf{K}_0\boldsymbol{\Phi}_j = \lambda_j\mathbf{M}_0\boldsymbol{\Phi}_j$,
$j=1,\ldots,n_\mathrm{modes}$'' but $n_\mathrm{modes}$ is never stated
as a parameter. Its value must be specified for each benchmark.

\item \textbf{Inconsistent void pseudo-material parameters.}
Section~4.1 (simply supported beam) uses $E_\mathrm{min}=10^{-9}E_0$;
Section~4.2 (clamped beam) uses $E_\mathrm{min}=10^{-6}E_0$. The
basis for this difference is not explained.

\end{enumerate}

%-----------------------------------------------------------------------
\section{Validation and Benchmarking Weaknesses}
%-----------------------------------------------------------------------

\begin{enumerate}[leftmargin=*, label=\textbf{V\arabic*.}]

\item \textbf{Table~1 lacks statistical uncertainty.}
Runtime and memory values are labeled ``averaged over 10 independent
runs'' without standard deviations. At the $160\times20$ mesh, the
proposed method (5.1~s) is marginally slower than the Yuksel code
(4.7~s). Without variance information, this difference and others in
Table~1 cannot be interpreted statistically. Standard deviations (or
at minimum, ranges) must be reported; the source of run-to-run
variability in a deterministic MATLAB code should be explained.

\item \textbf{No eigenfrequency values across mesh resolutions.}
Table~1 reports runtime and iterations at four mesh sizes but does not
include the achieved $\omega_1$ for any method at any resolution.
The 8.6\% frequency gap is stated only for the $400\times50$ case.
A column for $\omega_1$ at each mesh size should be added so that the
mesh-dependence of the frequency gap can be assessed.

\item \textbf{No mesh independence for multi-mode benchmarks.}
The clamped-beam ($400\times50$) and building ($80\times240$) examples
are each tested at a single mesh. MAC-tracked gains, mode-ordering
shifts, and the $\alpha$-behavior conclusions should be verified at
at least one additional discretization.

\item \textbf{Low MAC$_2$ values in Table~5 not acknowledged.}
Table~5 reports $\mathrm{MAC}_2=0.520$ at $\alpha=0.25$ and
$\mathrm{MAC}_2=0.635$ at $\alpha=0.5$ for the building example.
These values indicate that the mode identified as best matching
$\boldsymbol{\Phi}_2$ shares only about 27\% and 40\% of its modal
energy with the reference, respectively. The frequency gains reported
for these configurations (3.32$\times$ and 3.09$\times$) are based on
poorly correlated modes. The authors should define a MAC validity
threshold and qualify or exclude gain figures below it.

\item \textbf{Initial clamped-beam eigenfrequencies not stated.}
Section~4.3 explicitly states $\omega_1^{(0)}=19.84$~rad/s and
$\omega_2^{(0)}=90.93$~rad/s for the building domain. The analogous
values for the clamped-beam domain are absent from Section~4.2.
Table~3 gain factors cannot be independently verified without them.

\item \textbf{Headline claims must be scoped.}
\begin{itemize}
  \item ``within 8.6\% of the bound-formulation optimum'' should be
    qualified to the simply supported beam benchmark at $400\times50$.
  \item ``completing the optimization 7.1$\times$ faster'' is specific to
    $400\times50$; at $160\times20$ the proposed method is marginally
    slower than the Yuksel code.
  \item ``reaching up to 4.61$\times$ over the initial solid-domain
    frequency'' refers to the MAC-tracked $\Phi_1$-type gain, not the
    conventional sorted first eigenfrequency. This distinction should
    appear in the abstract.
\end{itemize}

\end{enumerate}

%-----------------------------------------------------------------------
\section{Mathematical Notation and Rigor}
%-----------------------------------------------------------------------

\begin{enumerate}[leftmargin=*, label=\textbf{M\arabic*.}]

\item \textbf{Eigenvector notation inconsistency.}
The manuscript uses uppercase $\boldsymbol{\Phi}$ (Eqs.~5, 7,
Section~3.6) and lowercase $\boldsymbol{\phi}$ (Eqs.~2, 3, MAC
definition Eq.~9) for eigenvectors without a stated convention.
Equation~(9) uses both $\boldsymbol{\phi}_i$ and $\boldsymbol{\Phi}_j$
in the same expression, suggesting an intentional distinction (current
vs.\ reference eigenvector), but this is never defined. A convention
should be stated at first use and applied consistently.

\item \textbf{Sensitivity filter notation (Eq.~8).}
The filtered sensitivity symbol is introduced without a formal
definition. The denominator summation domain ($\sum H_{ei}$) should
state explicitly that the sum runs over the neighborhood $N_e$ of
element~$e$.

\item \textbf{Cross-table discrepancy at $\alpha=1.0$ (clamped beam).}
Table~2 gives $\omega_2=365.78$~rad/s; Table~3 gives the
$\Phi_1$-tracked mode (MAC$=0.995$) at topo-mode~3 with
$\widetilde{\omega}_1=368.2$~rad/s. The 2.4~rad/s difference
($\approx0.7\%$) between these quantities should be explained
(rounding, independent runs, or precision issue).

\item \textbf{Scaling claim not supported by reported data.}
Section~5.2 states ``both methods exhibit an empirical per-iteration
scaling of roughly $\mathcal{O}(n_e^{1.3})$.'' The per-iteration times
in Table~1 are consistent with exponents of approximately 1.06 (Olhoff)
and 0.8 (Yuksel) over the four reported mesh sizes, substantially
below~1.3. The claim should be supported with a log-log scaling figure
or the specific exponent values should be corrected.

\end{enumerate}

%-----------------------------------------------------------------------
\section{Reproducibility and Transparency}
%-----------------------------------------------------------------------

\begin{enumerate}[leftmargin=*, label=\textbf{R\arabic*.}]

\item \textbf{No versioned code archive.}
The GitHub URL is provided without a commit hash or archived release
DOI. Repository contents may change after publication. An immutable
archive (e.g., Zenodo DOI or GitHub release tag) should be provided.

\item \textbf{Hardware and environment not specified.}
CPU model, clock speed, RAM, operating system, MATLAB version, and
thread count are absent. These are necessary to assess the reported
timing figures.

\item \textbf{No script-to-figure mapping.}
The manuscript does not specify which script or configuration file
reproduces each figure or table entry.

\item \textbf{``No eigensolve machinery'' is inaccurate.}
The abstract states the method requires ``only a single eigensolve at
initialization.'' Algorithm~3.5 Step~5 performs an additional eigensolve
at convergence. The correct description is ``no eigensolve \emph{inside}
the optimization loop.''

\end{enumerate}

%-----------------------------------------------------------------------
\section{Required Revisions}
%-----------------------------------------------------------------------

\subsection*{Critical (must be resolved before reconsideration)}

\begin{enumerate}[leftmargin=*, label=\textbf{CR\arabic*.}]

\item Verify and correct Table~3 $\Phi_2$-tracked gain data ($\alpha=0.75$
  entry shows $1.73\times$, which is below the $\alpha=1.0$ entry of
  $1.99\times$). Provide diagnostic output (iteration history, MAC
  evolution) for the $\alpha=0.75$ clamped-beam case. Revise or qualify
  all monotonicity claims in the abstract, contribution list,
  Section~4.2, and conclusions to match the corrected or explained data.

\item Add a quantitative validation of the omitted load sensitivity
  (Eq.~6): compare the approximate gradient against the full
  design-dependent-load derivative, or against a periodically
  updated-mode variant, on at least one benchmark.

\item Relabel the Olhoff comparator as an in-house Olhoff-inspired
  implementation throughout the manuscript; provide a comparison table
  showing what each method includes or excludes; state quantitatively
  how the three deviations (Section~5.1) contribute to the reported
  speedup and frequency gap.

\end{enumerate}

\subsection*{Major (must be addressed)}

\begin{enumerate}[leftmargin=*, label=\textbf{MR\arabic*.}]

\item Add standard deviations to all runtime and memory entries in
  Table~1; explain the source of variability across the 10~runs; specify
  CPU, RAM, OS, MATLAB version, and thread count.

\item Add eigenfrequency ($\omega_1$) values for all three methods at all
  four mesh resolutions in Table~1 or a supplementary table.

\item State initial clamped-beam eigenfrequencies $\omega_1^{(0)}$ and
  $\omega_2^{(0)}$ in Section~4.2.

\item Define and apply a MAC validity threshold; qualify or exclude frequency
  gains for configurations below the threshold (particularly
  $\alpha=0.25$ and $\alpha=0.5$ in Table~5).

\item Provide a mesh convergence study (minimum two mesh sizes) for the
  clamped-beam or building multi-mode benchmarks.

\item Provide a frozen-mode reliability diagnostic: at minimum, one
  benchmark or analysis characterizing when the initial mode shape is
  a poor proxy for the optimal mode shape, or demonstrating that the
  three tested problem classes lie within the reliable regime.

\item Scope all headline quantitative claims to the specific benchmark,
  mesh, and metric to which they apply.

\end{enumerate}

\subsection*{Minor (should be corrected)}

\begin{enumerate}[leftmargin=*, label=\textbf{mn\arabic*.}]

\item Correct figure references in Section~4.1 (``Figs.~5--7'' should
  be ``Figs.~1--3'').
\item Define the eigenvector notation convention ($\boldsymbol{\Phi}$
  vs.\ $\boldsymbol{\phi}$) at first use and apply consistently.
\item Complete the Olhoff--Du (2014) reference (publisher location listed
  as ``???'').
\item Archive the code at a permanent DOI and report the commit hash.
\item Specify hardware and MATLAB configuration in Section~5 or a
  reproducibility statement.
\item Provide a script-to-figure mapping in the repository README.
\item Address the OC-vs-MMA optimizer switch: state which optimizer is
  used for the proposed method in each benchmark and justify the choice.
\item Correct or support the $\mathcal{O}(n_e^{1.3})$ scaling claim
  (Section~5.2) with a log-log plot or revised exponent values.
\item Justify or unify the void pseudo-material parameter differences
  ($E_\mathrm{min}=10^{-9}E_0$ vs.\ $10^{-6}E_0$) across benchmarks.
\item Correct the inaccurate ``no eigensolve machinery'' claim to
  ``no eigensolve inside the optimization loop.''
\item Clarify the 2.4~rad/s discrepancy between Tables~2 and~3 at
  $\alpha=1.0$ (clamped beam, $\omega_2=365.78$ vs.\
  $\widetilde{\omega}_1=368.2$~rad/s).

\end{enumerate}

%-----------------------------------------------------------------------
\section*{Cross-Review Note}
%-----------------------------------------------------------------------

This report is based on an independent dual-model review (GPT and Claude
Sonnet), with both models reaching MAJOR REVISION with HIGH confidence.
The decision and trusted issue list are derived from the consensus of
both models; issues raised by only one model are retained where they
are supported by manuscript-visible evidence and are marked accordingly
in the editorial record. No hallucinated or invalid reviewer comments
were identified in either model's pipeline. The consensus weighted score
is approximately 2.80 on a 1--5 scale; the MAJOR REVISION recommendation
reflects an editorial override based on the correctable nature of the
critical issues.

\end{document}
